\documentclass[UTF8,zihao=-4]{ctexart}
\usepackage[a4paper,margin=2.4cm]{geometry}
\usepackage{booktabs}
\usepackage{array}
\usepackage{enumitem}
\setlist{nosep}
\emergencystretch=2.5em
\pagestyle{plain}
\title{\bfseries 全国大学生数学建模竞赛\\[0.3em] AI 工具使用详情}
\author{参赛队伍}
\date{}
\begin{document}
\maketitle
\vspace{-2.5em}
\section*{一、所用 AI 工具的名称、版本或型号}
\begin{table}[h!]
\centering
\begin{tabular}{@{}p{4.2cm}p{4.6cm}p{6.4cm}@{}}
\toprule
工具名称 & 版本/型号 & 说明 \\
\midrule
Claude Code & CLI 命令行版（Anthropic） & 会话式 AI 编程助手，用于代码编写、调试与文本辅助 \\
DeepSeek & V4-flash 系列模型 & Claude Code 底层驱动模型 \\
Ollama & minicpm-v 视觉模型 & 本地部署，用于 PDF 渲染页面的视觉校验 \\
Python 3 & 3.11 + numpy/pandas/scikit-learn & 数据分析与建模环境（脚本由 AI 辅助编写，见支撑材料附录 C） \\
\bottomrule
\end{tabular}
\end{table}

\section*{二、具体使用目的和环节}
\begin{enumerate}[label=(\arabic*)]
\item 问题建模环节：围绕 A 题退化特征提取、多因子容量衰减模型、RUL 预测与梯次筛选编组的思路梳理与公式校对；
\item 代码编写与调试环节：支撑材料中 7 个 Python 脚本（\texttt{generate\_sim.py}、\texttt{q1\_solution.py}、\texttt{q2\_solution.py} 等）的框架搭建、报错修复与单元校验；
\item 数据处理环节：数据字段映射与格式规范（\texttt{convert\_to\_spec.py}）脚本生成，CSV 交接文件的一致性校验；
\item 论文排版环节：\LaTeX{} 文档结构、图表引用、附录程序代码嵌入的辅助调整；
\item 文献整理环节：参考文献要点摘要与引用编号整理。
\end{enumerate}

\section*{三、主要提示方式与使用过程说明}
以自然语言描述建模目标、数据接口与期望输出，要求 AI 给出可直接运行的代码并标注参数来源，对生成内容逐项询问依据。典型交互示例：
\begin{quote}\small
提示：\emph{“请按《数据格式规范 v1.1》将仿真宽表转换为 A$\to$B$\to$C 交接文件，dod 改为小数、新增 cond\_id 与 condition 工况串，并做 SOH/dod 范围断言校验。”}
\end{quote}
AI 据此生成 \texttt{convert\_to\_spec.py}，经人工审查字段映射无误后纳入支撑材料，作为问题 1--4 各成员的数据接口。

\section*{四、对 AI 输出的采纳、人工修改和核验的主要情况}
\begin{sloppypar}
\begin{enumerate}[label=(\arabic*)]
\item 核心建模思路、公式推导与参数选择均由参赛队成员主导，AI 仅承担辅助编写与润色工作；
\item AI 生成代码经人工逐行审查后采用，关键数值结果（如三因子组合权重 $\{T:0.29,\ C:0.39,\ D:0.32\}$、温度寿命 1218$\to$206 循环、容量再生幅度 18.9\% 等）均以 Python 独立复算核验，与论文正文数据一致；
\item 数据文件字段与取值范围经脚本断言校验（SOH、dod 等）通过；
\item 语言润色类输出直接采纳；其余 AI 输出均经人工核实后采用。
\end{enumerate}
\end{sloppypar}

\vfill
\begin{center}
\zihao{-5}
本文件依据《全国大学生数学建模竞赛 AI 工具使用规定（2026）》第四条编制，随支撑材料一并提交。
\end{center}
\end{document}
